\documentclass[11pt,a4paper]{report}

\usepackage{microtype}
\usepackage{xcolor}
\usepackage{subcaption}
\usepackage{amsmath}

% Editing macros — delete for the final draft
\newcommand{\fc}[1]{\textcolor{red}{[FactC: #1]}}
\newcommand{\ct}[1]{\textcolor{orange}{[CITE: #1]}}
\newcommand{\app}[1]{\textcolor{blue}{[APP: #1]}}
\newcommand{\sr}[1]{\textcolor{teal}{[Sreword: #1]}}
\newcommand{\fig}[1]{\textcolor{purple}{[FIG: #1]}}
\newcommand{\ai}[1]{\textcolor{brown}{[AI: #1]}}

\begin{document}
%=======================================================================================
\chapter{Related Work and Background} \label{chap:related_work_and_background}

\section{Forest Growth and Chapman richards equation}
- top height, age (being main driver), yield class and stand development;
- biotic vs abiotic factors effecting growth (cause cocnsitent changes to growth, or deviations, trees do recover however and jump back to normal growth, or completely stagnate and enter a different growth trajectory) 
- why lidar data could be useful 
- broad influences of climate, terrain and management (satellite data);
%NOTE: question: are the envirnomental data sources mostly lidar or others?
- Chapman–Richards curve and parameters; (sigmoidal growth, say the equation and what the parameters are ymax: asymptote; k: growth rate; p: shape.
- connection to yield class and site index curves? biological interpretation of its shape, volume connection 
- strengths of CR, biolically plausible general form, 
- Cons of this model, pooled curve hides local heterogeneity, parameter uncertainty and ymaxk trade off? 
- why its relevant 


%------------------------------------------------
\section{Data driven and physics guided}
- Previous machine-learning approaches to forest prediction.
- Evidence comparing ML with conventional models.
- Physics-informed and hybrid approaches.
- Whether physical constraints improve prediction or mainly improve consistency.
- Importance of comparing a PINN with a physics-disabled DNN.
- How this literature motivates your CR–DNN–PINN comparison.


%------------------------------------------------
\section{Spatial evaluation and growth attribution}
- Why spatial dependence can make random validation optimistic.
- Previous use of spatial validation and Moran’s I.
- Previous work relating growth variation to environmental and management factors.
- Local or geographically weighted approaches relevant to GNNWR.
- How this motivates your spatial evaluation and two attribution avenues.


%------------------------------------------------
\section{Research gap and study position}
A short synthesis:
- What previous studies establish.
- What remains uncertain.
- What your experiments test differently.




%=======================================================================================

\end{document}
